% ============================================================
%  ch09_future.tex — Future Prospects
%  Requires: macros.tex
%  Estimated length: 6–10 pages
%  Rewritten: VD-07W, 2026-03-07
% ============================================================

\chapter{Future Prospects}\label{ch:future}

The defect-variable framework identifies seven falsifiable
predictions (extended to eleven with the tilted-FLRW programme)
and maps each to a specific combination of next-generation
experiments.  The cancellation of CMB-S4 in July
2025 and the LiteBIRD launch target of FY2032 (early 2030s) have
reshaped the experimental programme, concentrating decisive
discriminating power in the 2027--2030 window when five
independent probes---Euclid, DESI, SPHEREx (launched March 2025;
catalog-release timeline pending), Vera Rubin, and the
Simons Observatory---will simultaneously constrain the intrinsic
dipole, the filling fraction, and the preferred-axis alignment.

The conditional structure of the evidence
(Section~\ref{sec:reframing}) generates two distinct experimental
programmes depending on the outcome.  If the anomaly is confirmed
by independent surveys at $\geq 5\sigma$, the predictions below
become quantitative tests of the tilted-Bianchi interpretation.
If the anomaly is refuted (CatWISE-specific significance reduced
below $2\sigma$ and radio surveys similarly degraded), the defect
variables revert to type-independent upper bounds on cosmic
anisotropy, and the framework's contribution is the
three-bound hierarchy, the constraint-space geometry, and the
$T_{\mathrm{eff}}$ formalism---results that hold regardless of
the anomaly's fate.

The native low-ell morphology atlas remains the boundary between
scalar compatibility diagnostics and family-level morphology claims.
The current repository can define adapter schemas, artifact lanes,
null/covariance gates, and tomographic forecasts, but it does not
produce native low-ell transfer functions or morphology-atlas
comparisons.  Future promotion requires native atlas entries,
matched masks, full covariance, external nulls, and family-equivalence
tests.


% ============================================================
\section{Testable predictions}\label{sec:seven-predictions}
% ============================================================

The post-FLRW interpretation
(Section~\ref{sec:reframing} of Chapter~\ref{ch:dipole})
generates eleven predictions (P1--P11) testable with current
or funded experiments.  The first seven (P1--P7) target the
geometry-frame anisotropy; the four new predictions
(P8--P11, from the Tsagas bridge) target the perturbative
tilted-FLRW programme.

\emph{P1 (Intrinsic dipole excess).}
$\eps_1^{(\mathrm{int})} > 0$: the matter number-count dipole
exceeds the kinematic prediction across independent surveys at
different wavelengths and redshifts.  Currently confirmed at
$5.4\sigma$ (B\"ohme et~al.~2025).  The prediction is that the
excess persists in Euclid, Vera Rubin, and SKA catalogues.

\emph{P2 (Acceleration content).}
$\etaud \approx 8.3\%$: the matter-frame four-acceleration
contributes a specific fraction of the observed CMB dipole.
Testing this prediction requires distinguishing the
acceleration dipole from the kinematic dipole through
frequency-dependent or polarisation-dependent signatures.
No funded experiment can test this prediction decisively before
$\sim$2040.

\emph{P3 (Filling fraction).}
$\FF \geq 5$--$30\%$: the tilt occupies a non-negligible
fraction of the MES-allowed anisotropy budget, with $\FF(z)$
evolving under the DCP dynamics.  Six large-scale structure
surveys will constrain $\FF$ through redshift-binned dipole
measurements.

\emph{P4 (Low-$\ell$ B-mode).}
The tilt generates a B-mode polarisation signal at $\ell < 10$
with amplitude $C_\ell^{BB} \sim 10^{-12}\;\mu\mathrm{K}^2$,
approximately $10^6$ times below current sensitivity.  LiteBIRD
may achieve $\sim 1\sigma$ sensitivity by $\sim$2037.

\emph{P5 (Growing tensor mode).}
The regular tensor growing mode in Bianchi~VII$_h$ produces a
distinctive quadrupole-dominated pattern at
$(\sigma/H) \sim 10^{-6}$.  The Simons Observatory may reach
$\sim 2\sigma$ by 2030.

\emph{P6 (Axis clustering).}
The clustering of anomaly directions within $\sim 30^\circ$ is
testable by any experiment that independently measures a
cosmological preferred axis.  Euclid, Vera Rubin, and SKA can
each achieve $> 3\sigma$ directional constraints.

\emph{P7 ($\beta(z)$ evolution).}
The redshift dependence of the tilt rapidity distinguishes
between the four $\beta(z)$ models of
Section~\ref{sec:filling-z}.  Redshift-binned dipole
measurements from Euclid, DESI, and SKA will discriminate
between these scenarios by $\sim$2029.

% === INLINED: ch09_tilted_predictions.tex ===
% ============================================================
% ch09_tilted_predictions.tex
% New predictions P8–P11 and expanded forecast matrix
% Insert into ch09_future.tex:
%   (a) P8–P11 after P7 (line 83)
%   (b) Expanded forecast text after existing forecast section (line 108)
% ============================================================


% ── P8–P11: Insert after P7 (before the forecast section) ──

\emph{P8 (Dipolar deceleration scaling).}
The apparent deceleration
$\Delta q = (\beta/9)(\lambda_H/d)^3$
(equation~\eqref{eq:disc-Dq}) predicts a power-law decay with
survey depth, $\Delta q \propto d^{-3}$.  The Colin et~al.\
exponential fit $e^{-z/S}$ with $S = 0.0262$ reproduces the same
behaviour within $20\%$ for $d \approx 200$--$400\;\mathrm{Mpc}$
but diverges below $d \lesssim 150\;\mathrm{Mpc}$.  Rubin LSST
will accumulate $> 10^4$ SNe~Ia per $\Delta z = 0.02$ bin in the
range $0.02 < z < 0.15$, yielding
$\delta(\Delta q) < 0.1$ per bin---sufficient to distinguish
the Tsagas power law from the Colin exponential at $> 3\sigma$
by $\sim\!2029$.  DESI DR3+, expected in 2027, will independently
constrain the isotropic $q_0$ monopole through BAO + SN
combination, providing the zero-tilt baseline against which the
dipolar $\Delta q$ is measured.

\emph{P9 (Peculiar Jeans scale).}
The peculiar Jeans length
$\lambda_J^{\mathrm{pec}} = \lambda_H(\beta/(9q))^{1/3}
\approx 280$--$400\;\mathrm{Mpc}$
(Table~\ref{tab:disc-tilted-H0}) defines the transition scale above
which all observers agree on $q$.  If $\lambda_J^{\mathrm{pec}}$
is physical, the dipolar $q$ anisotropy must vanish for
$d > \lambda_J^{\mathrm{pec}}$.  Euclid will observe
$\sim\!2000$ SNe~Ia at $z > 0.8$
($d > 2700\;\mathrm{Mpc} \gg \lambda_J^{\mathrm{pec}}$),
providing a ``clean zone'' where the tilt correction is
negligible ($\Delta q < 10^{-3}$).  If a dipolar $q$
persists at $d > 600\;\mathrm{Mpc}$, the Tsagas mechanism
alone cannot explain it, and an alternative source of
anisotropy must be invoked.

\emph{P10 (Directional $H_0$).}
The tilted Hubble rate
$\hat{H}(\hat{n}) = H[1 + (\beta/3)(\lambda_H/d)\cos\theta]$
(equation~\eqref{eq:disc-H-dir}) predicts a $\cos\theta$
dipole in the locally measured $H_0$, with amplitude
$\sim\!5\;\mathrm{km\,s}^{-1}\mathrm{Mpc}^{-1}$ at
$d \approx 50\;\mathrm{Mpc}$ and maximum toward the CF4
bulk-flow direction.  The eROSITA all-sky X-ray cluster survey
(in safe mode since 2022-02-26; DR1 released 2024-01-31, DR2
planned mid-2026) has already catalogued $\sim\!10^5$ clusters
whose directional $H_0$ analysis could improve the Migkas
et~al.\ (2021) detection ($9\%$ variation, $5.4\sigma$) to
$\sim\!10\sigma$ precision and test the predicted $d^{-1}$
decay---provided future data releases resume or the existing
catalog proves sufficient.  SKA continuum surveys will provide
a complementary radio-selected $H_0$ map at $z < 0.1$ with
independent systematics.

This prediction faces a mixed observational landscape.
Stiskalek, Desmond \& Lavaux~\cite{Stiskalek2026} find a
zero-point dipole of $4.1 \pm 0.9\%$ in the CF4~W1 sample
($877{:}1$ Bayesian odds favouring anisotropy), but show
that a radially varying velocity dipole accounts for the
signal---attributing it to local flow features rather than a
cosmological $H_0$ anisotropy.  This interpretation is
\emph{consistent} with the tilt prediction, which posits
precisely a local velocity-induced $d^{-1}$ dipole rather
than a constant cosmological offset.  Conversely,
Sah et~al.~\cite{Sah2025} report a $> 5\sigma$ dipolar
$H_0$ variation in Pantheon+ at $0.023 < z < 0.15$ toward
the CMB dipole direction, interpreting it as a
general-relativistic effect of the local bulk flow.  The
parameter space is therefore tightening from the
cosmological-anisotropy side (constant dipole disfavoured)
while remaining open---and potentially
supportive---on the local-flow side ($d^{-1}$ decay
consistent with tilt).
We classify P10 as \emph{parameter space refined}: the
predicted scaling ($d^{-1}$ decay) is the critical
discriminant, and no analysis has yet tested this specific
functional form against the alternatives.

\emph{P11 (Age-bias disentanglement).}
The three-term decomposition
(equation~\eqref{eq:disc-q-decomp}) predicts that the age-bias
and tilt contributions to the apparent deceleration are
separable through their different dependences on host-galaxy
age and sky direction.  The age bias produces an isotropic,
redshift-growing signal that vanishes in evolution-free
subsamples; the tilt produces a dipolar, depth-decaying signal
that persists regardless of host properties.
Son et~al.~\cite{Son2025,Son2025b} have established the
progenitor age bias at $5.5\sigma$ significance across $> 300$
host galaxies (Part~I) and report alignment with DESI BAO
results suggesting a non-accelerating component at
$> 9\sigma$ (Part~II).  The joint directional-plus-age
disentanglement awaits Rubin-scale data.
Rubin LSST will provide $> 20{,}000$ SN~Ia host galaxies with
measured stellar population ages and full-sky coverage by
$\sim\!2030$.  A bivariate fit in (host age, $\cos\theta$) can
separate the two effects at $> 3\sigma$ if either contributes
$\Delta q \gtrsim 0.2$.  This test is the sharpest discriminant
between the Yonsei age-bias programme (Route~A) and the Tsagas
tilted-observer programme (Route~B), and will determine whether
the two effects are additive, redundant, or mutually exclusive.


% ── Expanded forecast text: append after existing forecast paragraph ──

The four new predictions expand the forecast matrix
(Table~\ref{tab:forecast-expanded}).  Two experiments emerge as
``golden'' probes of the tilted-FLRW extension:
Rubin LSST tests P8, P9, P10, and P11 simultaneously (through
$z$-binned dipolar $q$, host-age cross-correlation, and local
$H_0$ mapping), while Euclid provides the high-$z$ clean zone
for P9 and the spectroscopic BAO anchor for P8.

\begin{table}[t]
\centering
\footnotesize
\setlength{\tabcolsep}{3pt}
\caption{Extended experiment forecast matrix.  The first seven
predictions (P1--P7) are from
Section~\ref{sec:seven-predictions}; P8--P11 are derived from
the tilted-FLRW programme
(Section~\ref{sec:disc-tsagas}).
Symbols: $\checkmark\checkmark$ = primary sensitivity
($> 5\sigma$ projected), $\checkmark$ = secondary
($2$--$5\sigma$), --- = no direct sensitivity.}
\label{tab:forecast-expanded}
\begin{tabular}{lccccccccccc}
\toprule
Experiment &
  P1 & P2 & P3 & P4 & P5 & P6 & P7 &
  P8 & P9 & P10 & P11 \\
\midrule
LiteBIRD &
  $\checkmark\checkmark$ & $\checkmark$ &
  --- & $\checkmark$ & $\checkmark\checkmark$ &
  $\checkmark$ & --- &
  --- & --- & --- & --- \\
Simons Obs.$^{b}$ &
  $\checkmark$ & --- &
  --- & --- & $\checkmark$ &
  $\checkmark$ & --- &
  --- & --- & --- & --- \\
Euclid &
  $\checkmark$ & --- &
  $\checkmark\checkmark$ & --- & --- &
  $\checkmark$ & $\checkmark\checkmark$ &
  $\checkmark$ & $\checkmark\checkmark$ &
  --- & $\checkmark$ \\
DESI DR3+ &
  --- & --- &
  $\checkmark$ & --- & --- &
  --- & $\checkmark$ &
  $\checkmark$ & $\checkmark$ &
  --- & --- \\
Vera Rubin &
  --- & --- &
  $\checkmark\checkmark$ & --- & --- &
  $\checkmark\checkmark$ & $\checkmark\checkmark$ &
  $\checkmark\checkmark$ & $\checkmark$ &
  $\checkmark$ & $\checkmark\checkmark$ \\
SKA &
  --- & --- &
  $\checkmark\checkmark$ & --- & --- &
  $\checkmark\checkmark$ & $\checkmark\checkmark$ &
  --- & --- &
  $\checkmark$ & --- \\
eROSITA &
  --- & --- &
  --- & --- & --- &
  $\checkmark$ & --- &
  --- & --- &
  ---$^{a}$ & --- \\
SPHEREx$^{c}$ &
  --- & --- &
  $\checkmark$ & --- & --- &
  $\checkmark$ & $\checkmark$ &
  --- & --- & --- & --- \\
\bottomrule
\end{tabular}
\\[4pt]
{\footnotesize $^{a}$\,eROSITA has been in safe mode since
2022-02-26; DR1 was released 2024-01-31, DR2 planned mid-2026.
P10 sensitivity depends on resumed survey operations and is
therefore downgraded from the original forecast.\\
$^{b}$\,Simons Observatory: LAT first light Feb 2025;
entering science operations 2025--2026; no published results
at the first-light milestone. Sensitivities are forecast-level.\\
$^{c}$\,SPHEREx launched 2025-03-11; all sensitivities are
catalog-release-dependent.}
\end{table}

The extended matrix reveals a clear division of labour.  CMB
experiments (LiteBIRD, Simons Observatory) remain the primary
probes of the geometry-frame predictions (P1, P2, P4, P5), where
the tilted-FLRW extension contributes nothing new.  Large-scale
structure surveys (Euclid, Vera Rubin, DESI, SKA) now test
\emph{eleven} predictions spanning both the homogeneous defect
identity (P3, P6, P7) and the perturbative tilted extension
(P8--P11).  The eROSITA cluster catalog, if augmented by future
data releases (see table note), could contribute to directional
$H_0$ (P10); in its absence, Vera Rubin and SKA serve as the
primary probes.

Vera Rubin is the single most powerful experiment for the
tilted-FLRW predictions, testing P8, P9, P10, and P11 at the
$\checkmark$ or $\checkmark\checkmark$ level.  The decisive
test---the bivariate (host age, direction) analysis of P11---requires
the combination of high-statistics SN photometry with host-galaxy
spectroscopy, which Rubin will provide at unprecedented scale.
The first results from Rubin's 5-year survey are expected by
$\sim\!2030$, coinciding with DESI~DR3 and Euclid~DR2, creating
a second transformational window (2029--2032) alongside the
CMB-focused 2033--2037 window (LiteBIRD).

% === END INLINED ===


% ============================================================
\section{Experiment-by-experiment assessment}
\label{sec:experiment-assess}
\label{sec:forecast-experiments}
% ============================================================

\subsection{Simons Observatory}

The Simons Observatory~\cite{SimonsObs2019} achieved
LAT first light in February 2025,
joining three Small Aperture Telescopes that have been observing
since mid-2024.  The observatory is entering full science
operations through 2025--2026; no published scientific results
had appeared at the first-light milestone.  Enhanced LAT
upgrades (doubled detector count) are expected by 2028.

For the defect framework, the Simons Observatory tests P5
(growing tensor mode) through improved quadrupole and octupole
polarisation measurements, and P4 (low-$\ell$ B-mode) at
reduced sensitivity relative to LiteBIRD.  The sensitivity to
the regular tensor growing mode reaches $\sim 2\sigma$ at
$(\sigma/H) \sim 5 \times 10^{-7}$, within the detection
window of Section~\ref{sec:growing}.


\subsection{LiteBIRD}

LiteBIRD~\cite{LiteBIRD2023} is a JAXA-led
satellite mission targeting full-sky CMB polarisation at
$2 \leq \ell \leq 200$ with $\delta r < 10^{-3}$ sensitivity.
The launch target is FY2032 (early 2030s); Key Decision Point~2
(KDP2) was passed in September 2025, confirming the mission
timeline.  The cancellation of
CMB-S4 further elevates LiteBIRD's role as the primary
space-based CMB polarisation mission.

For the defect framework, LiteBIRD is the single most powerful
CMB experiment.  It tests P1 (intrinsic dipole separation via
low-$\ell$ polarisation), P4 (B-mode from tilt), and P5
(growing tensor mode at $> 5\sigma$).  The improvement factor
over the Saadeh vorticity bound is projected at $\sim 100\times$,
pushing $\Wstd$ below $10^{-23}$ and further compressing the
VII$_h$ constraint volume through the Frobenius coupling.
Data release is expected $\sim$2036.


\subsection{Euclid}

Euclid~\cite{Euclid2024} is an ESA photometric and
spectroscopic survey covering $15{,}000\;\mathrm{deg}^2$ to
$z \sim 2$.  The Q1 intermediate data release
(2025-03-19) covered 53 deg$^2$ from early operations;
Q2 (expected 2026-06-24) will expand coverage substantially.
The full DR1 is scheduled for October 2026.

For the defect framework, Euclid tests P1 (number-count dipole
at $z \sim 1$--$2$, independent of CatWISE), P3 (filling
fraction via redshift-binned $\beta(z)$), P6 (directional
alignment with the CatWISE/radio dipole), and P7
($\beta(z)$ discrimination between the four evolution models).
The spectroscopic galaxy sample also enables a cross-check
of the intrinsic dipole excess (P1) at optical/near-infrared
wavelengths.
The projected dipole sensitivity ($\sigma_{\eps_1} \sim
0.1 \times 10^{-3}$) would reduce the $\FF$ uncertainty by a
factor of $\sim 3$ relative to the current CatWISE measurement.


\subsection{SKA Phase~1}

The Square Kilometre Array Phase~1
\cite{SKA2020} will
survey $\sim 10^8$--$10^9$ radio continuum sources over
$30{,}000\;\mathrm{deg}^2$.

For the defect framework, SKA is the definitive number-count
dipole experiment.  With hundreds of millions of sources, the
shot noise limit drops to $\sigma_d \sim 10^{-4}$, enabling
$> 10\sigma$ sensitivity to the kinematic-vs-intrinsic dipole
separation.  SKA tests P1, P3, P6, and P7.  Data release is
expected from $\sim$2029.


\subsection{DESI}

The Dark Energy Spectroscopic Instrument
\cite{DESI2024} is obtaining
spectroscopic redshifts for $\sim 40 \times 10^6$ galaxies and
quasars at $0.1 < z < 3.5$.  DR1 became public on
2025-03-19, comprising 18.7 million spectra from the first
year of the main survey; DR2 BAO results (first three years)
were released the same date, with full DR2 cosmology chains
following in October 2025.

For the defect framework, DESI provides the redshift
precision needed for P7 ($\beta(z)$ evolution) through
spectroscopic dipole measurements in narrow redshift bins.
The large quasar sample at $z > 2$ extends the dipole
measurement to the epoch where the DCP tilt instability
predicts the strongest signal.  We note that
von~Hausegger \& Dalang~\cite{vonHauseggerDalang2025}
have shown that redshift-selection effects can be substantial
in binned kinematic dipole measurements---in extreme cases
even reversing the apparent dipole direction---so
redshift-binned forecasts for P3 and P7 must incorporate
selection corrections before being interpreted
quantitatively.


\subsection{Vera Rubin LSST}

The Vera C.\ Rubin Observatory Legacy Survey of Space and
Time~\cite{Rubin2019} released Data Preview~1 (DP1) in
June 2025, with full survey operations beginning no earlier
than late 2025.  The survey will image
$\sim 2 \times 10^{10}$ galaxies over
$18{,}000\;\mathrm{deg}^2$ to $r \sim 27.5$.

For the defect framework, the photometric dipole
measurement tests P1 at optical wavelengths (independent of
radio and infrared).  The depth and sky coverage provide the
best directional constraint on the dipole axis (P6) and the
first photometric measurement of $\beta(z)$ (P7).


% ============================================================
\section{The 2027--2030 window}\label{sec:window}
% ============================================================

By 2030, five independent probes will have simultaneously
constrained $\eps_1$, $\FF(z)$, and directional alignment,
providing the multi-probe redundancy needed to distinguish a
physical signal from survey systematics.  If the dipole anomaly
is physical, the convergence of independent measurements at
different wavelengths and redshifts will elevate the significance
well beyond the current $> 5\sigma$.  If the anomaly is
systematic, the diversity of survey methods will expose the
source.  the conditioned legacy figure gallery in Appendix~\ref{app:conditioned-legacy-figure-gallery} summarises the
projected timeline for each experiment.

The four-dimensional constraint space
$(\Sigstd, \Wstd, \Astd, \beta)$ will be progressively reduced:
the vorticity dimension by LiteBIRD's low-$\ell$ B-mode
($\Wstd \to 10^{-23}$); the tilt dimension by Euclid and SKA
redshift-binned dipoles; the shear dimension by the growing
tensor mode search.  The acceleration dimension
$\Astd$ remains observationally inaccessible until a
PIXIE-class spectrophotometer is funded.


% ============================================================
\section{Open problems}\label{sec:open-problems}
% ============================================================

Seven problems remain open within the defect-variable framework.

First, the sole published numerical bound on the cosmological
four-acceleration ($|\dot{u}|/\Theta < 10^{-4}$) dates to the
COBE era~\cite{StoegerAG1997} and has never been updated with
WMAP or Planck data.  Updating this bound using the Planck PR4
maps would close the last gap in the MES three-bound hierarchy.

Second, the homogeneous Bianchi framework cannot constrain
inhomogeneous backreaction.  A systematic connection between the
departure parameter $x$ and the Buchert backreaction
$\mathcal{Q}_\mathcal{D}$---perhaps through a coarse-graining
procedure---remains an open theoretical problem.

Third, the direction-marginalised Bayes factor
$\ln\mathcal{B} \approx +26$ (VER05 primary) discards potentially
model-separating directional information.  A joint pixel-level
analysis using the AniCLASS $a_{\ell m}$ output could
significantly sharpen the model selection; AniCLASS is
publicly available now at
\texttt{github.com/CyrilPitrou/aniclass} and the BASS Phase~3A
oracle validation protocol is immediately executable.

Fourth, the $\beta$-gap between the maximum curvature-sourced
tilt ($\beta_{\max}^{(\mathrm{curv})} \sim 10^{-7}$) and the
anomaly ($\beta \sim 10^{-3}$) implies that standard Bianchi~V
physics cannot generate the observed tilt from spatial curvature
alone.  Identifying the physical mechanism is the central
theoretical challenge.

Fifth, the effective $\etaud$ for a realistic multi-component
universe is epoch-dependent and must be integrated over the
photon's line of sight---a computation that has not been
performed.

Sixth, a Bayesian reformulation of $\FF$ as a posterior
probability rather than a ratio of algebraic bounds would
provide a more principled statistical interpretation.

Seventh, the Boltzmann-solver landscape for nearly isotropic
Bianchi cosmologies has advanced with the publication of
AniLoS/AniCLASS~\cite{Vicente2026}, which implements
scalar--vector--tensor (SVT) perturbation modes for
Bianchi types VII$_h$, V, VII$_0$, and IX---achieving roughly
$10\times$ speedup over the original AniLoS integrator.
ABSolve, the code underlying the Saadeh
et~al.~\cite{Saadeh2016b} template analysis, has no public
repository.  Neither solver implements genuinely tilted-fluid
Bianchi degrees of freedom: AniLoS/AniCLASS operate in the
SVT vector-mode sector, which is physically distinct from the
fluid-tilt sector exploited in this work.  A future BASS
(Bianchi Anisotropy Summary Statistic) solver that couples
the defect-variable likelihood to a Boltzmann integrator
with fluid-tilt degrees of freedom would therefore occupy a
genuine novelty niche---extending to a regime that no
existing public code covers.

Eighth, the reduced line-of-sight (R-LOS-T0) surrogate developed
in this work is structurally specified but numerically dormant:
the conditioned legacy figure gallery in Appendix~\ref{app:conditioned-legacy-figure-gallery} shows that the channel activates
only when $\Sigma^2 \gtrsim 10^{-11}$, which is 14 orders of
magnitude above the current posterior.
the conditioned legacy figure gallery in Appendix~\ref{app:conditioned-legacy-figure-gallery} confirms that the
projector--LOS divergence is 16--98\% in the stress-test regime
with kernel robustness CV $= 5.5\%$.  The on/off evidence rerun
produces $|\Delta\ln\mathcal{B}| = 0.10$---negligible.  The
module is ready; the data do not yet require it.

Ninth, the S-1M-SHADOW scalar spectral extension
(the conditioned legacy figure gallery in Appendix~\ref{app:conditioned-legacy-figure-gallery}) demonstrates that a bounded
one-mode $y$-distortion basis has a genuine admissible wedge
($\kappa \in (-1.0, 16.7)$, width $17.7$), exact $T_\mathrm{eff}$
recovery at $\varepsilon_{\mathrm{spec}} = 0$, and 13.5\%
history-sensitive divergence across $\Sigma^2$ values.  The shadow
is degenerate with $T_\ell$ at a fixed history but distinguishes
different source histories---the same physics as the R-LOS-T0
result.  Both extensions point to a common conclusion:
source-history timing information is structurally available within
the scalar sector but numerically dormant at current data precision.

Tenth, independent depth tomography from Quaia quasar samples
(the conditioned legacy figure gallery in Appendix~\ref{app:conditioned-legacy-figure-gallery}) can test whether the
tilt-rapidity detection persists across redshift shells.  If
$\beta$ is cosmological, the dipole amplitude should not fade
beyond the local-structure horizon ($z \gtrsim 0.1$).  This test
is deferred to the era of larger bulk-flow catalogues (e.g.\ CF4++)
but requires no new theory.


% === INLINED: ch09_solver_programme.tex ===
% ============================================================
% ch09_solver_programme.tex
% New future-directions sections for ch09
% Part of: Tetrad-Based Departure Decomposition (Jiwon Park)
% Session S7 draft — 2026-03-30
% ============================================================


\section{The anisotropic Boltzmann solver programme}
\label{sec:solver-programme-ch9}

The Phase~1.0 solver of Chapter~\ref{ch:results}
demonstrates the feasibility of a species-resolved
Bianchi Boltzmann calculation at production speed
($420$~ms per evaluation).  We now describe the
programme of extensions that transforms this prototype
into a full anisotropic Einstein--Boltzmann solver
capable of producing direction-dependent $a_{\ell m}$
spectra, self-consistent gravitational potentials, and
species-resolved transfer functions across all nine
Bianchi types.


\subsection{Competitive context}
\label{sec:competitive}

The field of anisotropic Boltzmann solvers as of
March~2026 is dominated by two codes.
AniCLASS~\cite{VicentePereiraPitrou2026}, published in
Physical Review D in February~2026, implements linearised
supercurvature-mode perturbations for Bianchi types
VII$_h$, V, VII$_0$, and IX, with species-resolved kinetic
hierarchies for photon temperature, polarisation, and
neutrinos.  ABSolve, the code underlying the
Saadeh~et~al.~\cite{Saadeh2016b} Planck template analysis,
supports types I, V, VII$_0$, and VII$_h$ but has no public
repository and has not been updated since~2016.  No other
code---including SONG, DISCO-EB, SymBoltz.jl, and
gCAMB---operates on anisotropic backgrounds.

Five capabilities of the planned solver have zero coverage
in any existing code:
exact (nonperturbative) Bianchi background evolution,
tilted-fluid Boltzmann transport,
a local peculiar-velocity patch formalism,
matrix line-of-sight integration for anisotropic
geometries, and per-species tangency diagnostics
$D_{s,\geq 2}$
(Section~\ref{sec:species-tangency}).  The
Maartens--Gebbie--Ellis~\cite{MaartensGE1999} nonlinear
hierarchy for tilted models has never been numerically
implemented.  The nearest emerging competitor is the
Oxford--Prague group (Mart{\'i}n, Skordis,
Ferreira~et~al.), whose tilted-dipole
analysis~\cite{MartinSkordis2025} uses Boltzmann-level
consistency checks and could evolve toward a full solver,
though their current framework does not solve the coupled
hierarchy.


\subsection{Phase~2.0: self-consistent Poisson}
\label{sec:phase20}

The Phase~1.0 solver uses an analytic transfer function
$\Psi(k, \eta)$ for the gravitational potential
(Section~\ref{sec:flrw-oracle}), which introduces a
$-8\%$ deviation from the CAMB reference value
(Section~\ref{sec:disc-oracle}).  Phase~2.0 replaces
this analytic prescription with a self-consistent
solution of the Poisson equation sourced by the
species-resolved anisotropic stress:
\begin{equation}\label{eq:poisson-sc}
k^2 \Psi = -4\pi G\, a^2
  \left[
  \sum_{s \in \mathcal{K}}
    (\rho_s + p_s)\, \Theta^{(s)}_0
  + \sum_{\alpha \in \mathcal{F}}
    \delta\rho_\alpha
  \right]\,,
\end{equation}
where the species contributions on the right-hand side
are evolved by the Bianchi hierarchy of
Section~\ref{sec:bianchi-quad}.  The target is
$|D_2^{\rm solver} - D_2^{\rm CAMB}| / D_2^{\rm CAMB}
< 2\%$, closing the oracle deviation.

Three additional deliverables accompany the
self-consistent Poisson step.
First, the full neutrino hierarchy with
$\ell_{\rm max,\nu} \geq 20$ replaces the current
two-moment approximation, capturing the free-streaming
transfer at all relevant multipoles.
Second, B-mode polarisation from the Bianchi tensor shear
extends the observable set from TT and EE to include BB,
enabling direct tests of the gravitational-wave
background predicted by anisotropic expansion.
Third, a nested sampling end-to-end integration test
couples the BianchiSolver to the dynesty sampler,
verifying that the evidence computation
$\ln \mathcal{B}$ is numerically stable under the
self-consistent potential.


\subsection{Phase~3.0: direction-dependent $a_{\ell m}$}
\label{sec:phase30}

The scalar-amplitude pipeline of VER06 compresses the
CMB anisotropy into a single number $D_2$ and marginalises
over direction.  Phase~3.0 lifts this compression by
computing the full $a_{\ell m}$ spectrum, which retains
the type-dependent angular pattern that distinguishes
Bianchi VII$_h$ (spirals) from Bianchi~I (pure
quadrupole) from Bianchi~IX (closed-universe hot spots).

The $a_{\ell m}$ computation requires the line-of-sight
integration of the species-resolved source functions
along each sky direction.  For a Bianchi model with
symmetry axis $\hat{n}^a$, the temperature anisotropy
along direction $\hat{e}^a$ is
\begin{equation}\label{eq:alm-los}
\frac{\Delta T(\hat{e})}{T_0}
= \int_0^{\eta_0} d\eta\,
  g(\eta)\,\big[
  \Theta^{(\gamma)}_0
  + \Psi
  + \hat{e}^a\, \Theta^{(\gamma)}_a
  + \hat{e}^{\langle a}\hat{e}^{b\rangle}\,
    \Theta^{(\gamma)}_{ab}
  + \cdots
  \big]
+ \text{(ISW)}\,,
\end{equation}
where $g(\eta) = \dot{\tau}\, e^{-\tau}$ is the
visibility function and the angular dependence of the
PSTF multipoles $\Theta^{(\gamma)}_{A_\ell}$ on the
symmetry axis $\hat{n}^a$ produces the type-dependent
$a_{\ell m}$ pattern.

AniCLASS provides the primary validation oracle for
this phase.  The cross-validation protocol of
Section~\ref{sec:aniclass-validation} extends from the
current $15$ grid points to a full type-by-type
comparison of $a_{\ell m}$ spectra at matched
cosmological parameters.

The $a_{\ell m}$ output enables a pixel-level template
fit to Planck and future CMB data, replacing the
direction-marginalised evidence with a
direction-inclusive Bayes factor.  This analysis could
break the equivalence classes identified in
Chapter~\ref{ch:robustness}---specifically, the
VII$_h$/V degeneracy in the scalar-amplitude pipeline---by
exploiting the distinct angular morphologies of the
different Bianchi types.


\subsection{Differentiation from AniCLASS}
\label{sec:differentiation}

AniCLASS is the nearest published competitor to the
planned solver.  The two codes share a common goal
(computing CMB anisotropy in Bianchi cosmologies) but
differ in five respects that define the planned solver's
unique contribution.

\emph{Background treatment.}  AniCLASS operates in the
linearised regime ($\sigma/H \ll 1$), exploiting the
mathematical correspondence between Bianchi modes and
FLRW supercurvature perturbations.  The planned solver
evolves the exact Bianchi background (scale factor,
shear, curvature) without a small-anisotropy expansion,
enabling access to the nonperturbative regime
$\sigma/H \gtrsim 10^{-3}$ where the linearised
correspondence breaks down.

\emph{Tilted-fluid dynamics.}  AniCLASS treats only the
vector-mode sector of the Bianchi perturbations, which is
physically distinct from the fluid-tilt sector.  The
planned solver implements the species-resolved tilt
through the dipole $\Theta^{(s)}_a$
(Section~\ref{sec:species-teff}), which couples to the
matter frame through the Thomson interaction
(Section~\ref{sec:thomson-coupling}) and enters the
departure identity through $\Otilt$.

\emph{Species-resolved $\Teff$.}  Both codes use
species-resolved kinetic hierarchies.  The planned solver
also provides the two-field neutrino treatment
($\Theta_\nu, \eta_\nu$;
Section~\ref{sec:species-teff}), the tangency diagnostic
$D_{s,\geq 2}$
(Section~\ref{sec:species-tangency}), and the connection
to the departure parameter pipeline through the
calibrated transfer function $T_2$
(Section~\ref{sec:species-D2}).

\emph{Local-patch formalism.}  AniCLASS computes the
global Bianchi $a_{\ell m}$ without accounting for the
observer's local peculiar velocity.  The planned solver
includes the bounded local-patch source
(Section~\ref{sec:disc-local-patch}), enabling the
TT-versus-EE separation between local and cosmological
contributions to the observed dipole.

\emph{Departure-parameter integration.}  The planned
solver feeds directly into the BASS/HTT/MIO inference
pipeline (Chapter~\ref{ch:pipeline}), computing
$\ln \mathcal{B}$ as a by-product of each evaluation.
AniCLASS can in principle be wrapped in an external
sampling framework, but no such integration exists.


\subsection{Scalar modes and Bianchi~I}
\label{sec:scalar-modes}

AniCLASS implements tensor ($|m| = 2$) and vector
($|m| = 1$) modes but not scalar ($m = 0$) modes.
The planned solver's exact-background approach does not
decompose the anisotropy into SVT modes; instead, it
evolves the full PSTF hierarchy coupled to the Einstein
equations, which automatically includes all mode types.
For Bianchi~I, which is excluded from AniCLASS because
all supercurvature wavenumbers $\nu_m = 0$ in the
linearised formalism, the planned solver produces
non-trivial output because the exact background shear
$\sigma_{ab}$ sources the hierarchy directly.  The
Phase~1.0 validation (Section~\ref{sec:disc-solver})
confirms this: $D_2(\mathrm{BI}, \Sigma^2 = 10^{-8})
= 5.50 \times 10^{-5}$ in hierarchy$^2$ units.


% ============================================================
\section{Backreaction and the beyond-Bianchi programme}
\label{sec:backreaction-programme}
% ============================================================

The Bianchi framework is spatially homogeneous by
construction: it constrains the global departure from
FLRW but cannot address the spatially inhomogeneous
backreaction that arises from the nonlinear coupling of
sub-Hubble perturbations.  We now outline the programme
that connects the homogeneous departure decomposition
to the Buchert spatial-averaging framework and ultimately
to a backreaction solver.


\subsection{The Buchert--Bianchi dictionary}
\label{sec:buchert-bianchi}

The Buchert equations (Section~\ref{sec:buchert-equations})
relate the domain-averaged expansion rate
$\langle \Theta \rangle_{\mathcal{D}}$ to the kinematic
backreaction $\mathcal{Q}_{\mathcal{D}}$ and the averaged
spatial curvature
$\langle \mathcal{R} \rangle_{\mathcal{D}}$ through the
integrability condition
\begin{equation}\label{eq:buchert-integrability}
\partial_t \mathcal{Q}_{\mathcal{D}}
+ 6\,H_{\mathcal{D}}\, \mathcal{Q}_{\mathcal{D}}
+ \partial_t
  \langle \mathcal{R} \rangle_{\mathcal{D}}
+ 2\,H_{\mathcal{D}}\,
  \langle \mathcal{R} \rangle_{\mathcal{D}}
= 0\,.
\end{equation}
The master departure identity of
Section~\ref{sec:master-id} provides a complementary
decomposition: $x_C = \Sigstd - \Wstd + \Otilt +
\Okaniso$ tracks the departure from FLRW in the
geometry-frame kinematic variables.  A natural
question is whether these two decompositions can be
connected by a dictionary.

The tilt sector of the departure identity maps to the
stress-energy backreaction $\mathcal{T}_{\mathcal{D}}$
in the generalised Buchert framework of
Buchert, Mourier \&
Roy~\cite{BuchertMourierRoy2020}, where
the extension to tilted fluids with vorticity introduces
additional backreaction terms sourced by the energy flux
$q_a$ and the tilt rapidity $\beta$.  The shear mapping
$\Omega_{\mathcal{Q}} \leftrightarrow \Sigstd$---where
$\Omega_{\mathcal{Q}}$ is the kinematic backreaction
density parameter---is structurally exact and algebraic
for homogeneous models.  For inhomogeneous models, the
correspondence becomes approximate, limited by the
mismatch between the Bianchi trace-free spatial curvature
${}^3\!S_{ij}$ and the Buchert scalar average
$\langle \mathcal{R} \rangle_{\mathcal{D}}$.

\begin{remark}[Structural mismatch]
\label{rem:bianchi-buchert-mismatch}
The Bianchi departure parameter $\Okaniso$ is a tensorial
quantity (the trace-free part of the spatial Ricci
tensor), while the Buchert averaged curvature
$\langle \mathcal{R} \rangle_{\mathcal{D}}$ is a scalar.
A complete dictionary would require Heinesen's
multipole cosmography
framework~\cite{Heinesen2021}, which decomposes
observables in a covariant multipole expansion without
assuming spatial homogeneity.  This bridge is deferred
to future work (Section~\ref{sec:open-problems}).
\end{remark}


\subsection{Domain-averaged effective sources}
\label{sec:domain-sources}

The species-resolved framework of
Section~\ref{sec:multi-fluid} provides the building
blocks for a backreaction calculation.  On a spatial
domain $\mathcal{D}$, the coarse-grained kinetic
variables are
\begin{equation}\label{eq:coarse-variables}
\begin{aligned}
\langle q_a \rangle_{\mathcal{D}}
&= \sum_{s \in \mathcal{K}}
  \langle q_s^a \rangle_{\mathcal{D}}
+ \sum_{\alpha \in \mathcal{F}}
  \langle q_\alpha^a \rangle_{\mathcal{D}}\,,\\[4pt]
\langle \pi_{ab} \rangle_{\mathcal{D}}
&= \sum_{s \in \mathcal{K}}
  \langle \pi_s^{ab} \rangle_{\mathcal{D}}
+ \sum_{\alpha \in \mathcal{F}}
  \langle \pi_\alpha^{ab} \rangle_{\mathcal{D}}\,,
\end{aligned}
\end{equation}
where each species contributes through its reconstructed
stress-energy (Section~\ref{sec:macro-recon}).  The
quadratic backreaction terms---such as the tilt-squared
effective stress
$\langle V_{\langle a} V_{b\rangle} \rangle_{\mathcal{D}}$
(Eq.~\ref{eq:tilt-squared})---arise naturally from the
second-order source library of
Section~\ref{sec:quadratic-sources}.

The backreaction solver would compute these
domain-averaged sources from the output of the Bianchi
hierarchy and feed them into the Buchert integrability
condition~\eqref{eq:buchert-integrability}, producing a
time-dependent backreaction history
$\mathcal{Q}_{\mathcal{D}}(t)$ that can be compared with
the departure parameter $x_C(t)$.  The key question is
whether the departure parameter, measured at the present
epoch, is consistent with the cumulative backreaction
integrated over the photon's line of sight.


\subsection{Error measurement hierarchy}
\label{sec:error-hierarchy-preview}

The accuracy of the solver at each phase is quantified by
a five-level error measurement hierarchy.  We state the
levels here and defer the detailed construction to
Chapter~\ref{ch:error-hierarchy}.

The first level, the \emph{state-space error}
$\varepsilon_{\rm state}$, measures the deviation of the
solver's PSTF multipoles from a reference solution (CAMB
in the FLRW limit, AniCLASS in the Bianchi regime).  The
second level, the \emph{source-translation error}
$\varepsilon_{\rm Th}$, measures the accuracy of the
$\Teff$-to-Thomson bridge in the photon sector
(Eq.~\ref{eq:exact-bridge}).  The third level, the
\emph{adequacy certificate} $\Delta_{\rm obs}^{\rm bound}$,
is the observable-level bound on the representation error,
computed from the tangency diagnostic $D_{s,\geq 2}$
and the Jacobian condition number $\kappa(J_s)$
(Section~\ref{sec:species-tangency}).  The fourth level,
the \emph{ladder-specific adequacy} $\Delta_q$ and
$\Delta_{N/E}$, quantifies the error incurred by the
adaptive $\ell$-band routing
(Section~\ref{sec:ell-band}).  The fifth level, the
\emph{observable error} $\varepsilon_{a_{\ell m}}$ and
$\varepsilon_{C_\ell}$, measures the deviation of the
final CMB power spectrum from the reference.

At each phase of the solver programme, the error budget
is dominated by a different level: Phase~1.0 is dominated
by the state-space error (the $-8\%$ oracle deviation);
Phase~2.0 targets the source-translation error (closing
the Poisson gap); Phase~3.0 targets the observable error
(matching AniCLASS $a_{\ell m}$ to $< 1\%$).


\subsection{Experimental alignment}
\label{sec:experimental-alignment}

The solver programme is designed to deliver results on a
timeline matched to the observational programme.  We
summarise the alignment between solver phases and
experimental milestones.

\emph{Euclid~DR1} (scheduled October~2026) will provide
$\sim\!2100~\mathrm{deg}^2$ of wide-survey photometric
data, enabling the first independent large-area
Ellis--Baldwin dipole test at optical/near-infrared
wavelengths.  Phase~2.0 targets completion before
Euclid~DR1, providing a self-consistent $D_2$ prediction
that can be directly compared with the Euclid dipole
measurement.

\emph{Simons Observatory} first-light data
($\sim\!2027$) will deliver high-resolution ground-based
CMB polarisation maps.  Phase~3.0 targets completion
before SO data release, enabling template fitting of
Bianchi $a_{\ell m}$ patterns against the SO EE and BB
spectra.

\emph{DESI} peculiar-velocity catalogues (DR2 and
beyond, $\sim\!2027$--$2028$) will extend the
CosmicFlows-4 bulk-flow measurement to larger volumes
and higher precision.  The species-resolved solver can
compute the predicted peculiar-velocity field from the
tilt sector $\Otilt$, providing a direct cross-check
between the CMB-derived and galaxy-derived constraints on
the departure parameter.

\emph{LiteBIRD} ($\sim\!2033$) will perform a full-sky
CMB polarisation survey with sensitivity to
$\Sigstd \sim 10^{-6}$, an order of magnitude beyond
current Planck constraints.  Phases~4.0--5.0 target
completion before LiteBIRD data release, enabling the
full adaptive-ladder and backreaction solver to be applied
to the highest-precision anisotropy data.

\emph{SKA Phase~1} ($\sim\!2028$--$2030$) will survey
the radio continuum dipole with sensitivity sufficient to
confirm or rule out the NVSS/RACS/LoTSS excess at
$> 10\sigma$.  The planned solver's direction-dependent
$a_{\ell m}$ output can be directly compared with the
SKA dipole direction and amplitude.

The critical timeline constraint is the window between
Euclid~DR1 (October~2026) and LiteBIRD ($\sim\!2033$):
within this interval, three or more of the seven
predictions of Section~\ref{sec:seven-predictions} will be
testable, and the solver must be ready to confront these
data with species-resolved, direction-dependent
predictions.

% === END INLINED ===

% === INLINED: figures_ch09.tex ===
% Figures for Chapter 9
% Updated: TF-D06, 2026-03-12
% Replaced: fig_experiment_timeline (P8-P11 added, eROSITA added)

\begin{table}[t]
\centering
\caption{Experiment forecast: earliest threshold crossings by prediction.
Predictions P8--P11 are new (TF-series).}
\label{tab:forecast}
\footnotesize
\begin{tabular}{llll}
\toprule
Prediction & Earliest threshold crossing & Experiment & Significance \\
\midrule
P1: Intrinsic dipole & 2028 & Simons Obs. & $3$--$5\sigma$ \\
P2: Accel.\ content & 2043 & PIXIE-class & $>3\sigma$ \\
P3: Filling fraction & 2027 & DESI & $3$--$5\sigma$ \\
P4: Tilt B-mode & 2037 & LiteBIRD & $<1\sigma$ \\
P5: Growing tensor & 2030 & Simons Obs. & $\sim 2\sigma$ \\
P6: Axis clustering & 2027 & DESI & $\sim 2\sigma$ \\
P7: $\beta(z)$ evolution & 2027 & DESI & $\sim 3\sigma$ \\
\midrule
P8: $\Delta q \propto d^{-3}$ & 2029 & Vera Rubin & $3$--$5\sigma$ \\
P9: Peculiar Jeans scale & 2030 & Vera Rubin & $\sim 3\sigma$ \\
P10: Directional $H_0$ & 2029 & Vera Rubin & $\sim 3\sigma$ \\
P11: Age-bias disentanglement & 2030 & Vera Rubin & $\sim 3\sigma$ \\
\bottomrule
\end{tabular}
\end{table}

% === ADDED 2026-03-22: R-LOS-T0 and S-1M-SHADOW future-work figures ===


% === ADDED 2026-04-24: research-plan extensions (F119, F120) ===

\subsection{Forecast Fisher matrix and tension-resolution timeline}
\label{sec:forecast-fisher-and-timeline}

Section~\ref{sec:forecast-experiments} mapped each Bianchi
prediction to the earliest forecast threshold-crossing experiment.
Two complementary
visualisations consolidate that information.

the conditioned legacy figure gallery in Appendix~\ref{app:conditioned-legacy-figure-gallery} compares forecast
Fisher 1$\sigma$ ellipses in the $(\beta, \sigma_*)$ plane, where
$\sigma_* \equiv \log_{10}(\sqrt{\Sigma^2_{\mathrm{std}}}/H)$ is the
shear-amplitude proxy used in the Planck-data $\Sigma^{2}$
posterior.  The ``current'' contour combines Planck~PR3, CatWISE and
CF4; subsequent contours show the Fisher contraction expected from
Euclid~DR1 (2027), Simons Observatory (2027), CMB-S4 and Vera~Rubin
LSST (2030), and LiteBIRD (2032).  The fiducial point is the VER05
posterior median $(\beta, \sigma_*) = (1.334 \times 10^{-3}, -5.0)$.
Even the most pessimistic (``current'') ellipse already separates the
two literature values $\beta_{\mathrm{CF4}} = 1.334 \times 10^{-3}$
and $\beta_{\mathrm{CMB}} = 1.233 \times 10^{-3}$ at $> 1\sigma$;
Euclid~DR1 separates them at $\sim 2.4\sigma$, and the CMB-S4
$+$~Rubin combination at $\sim 5\sigma$.

the conditioned legacy figure gallery in Appendix~\ref{app:conditioned-legacy-figure-gallery} re-organises the same
information by texture (alternative explanations of the dipole
anomaly) rather than by survey.  Each row corresponds to one
candidate texture (tilt-only, shear, shear$+$vortex, vortex-grow,
isocurvature defect, Khronon).  Coloured markers indicate the year at
which each survey is expected to push the relevant Bayes factor
beyond the historical Jeffreys high-support threshold
$|\ln\mathcal{B}| \ge 5$.  The tilt-only texture has already crossed
that forecast threshold (2024 -- Planck PR4 + CatWISE); shear textures
cross it in 2026--2027 with CF4++ and Euclid; the vortex-grow and
Khronon textures, which both rely on small-$\ell$ B-mode amplitude,
only cross it with LiteBIRD in 2032.

% === END INLINED ===
